\section{The Fundamental Theorem of Calculus}

\begin{callout}
  We take the following theorem to be given (as the proof has been
  \href{https://www.youtube.com/watch?v=dgfnO4uDR58}{provided}).
  \begin{theorem}[Fundamental Theorem of Calculus]
  \label{thm:ftc}
    \leavevmode
    \begin{enumerate}[label=(\roman*)]
      \item \label{thm:ftc-i} If $f : [a, b] \to \R$ is integrable, and $F : [a, b] \to \R$ satisfies $F'(x) = f(x)$ for all
        $x \in [a, b]$, then
        \[
          \int_{a}^{b} f = F(b) - F(a).
        \]

      \item \label{thm:ftc-ii} Let $g : [a, b] \to \R$ be integrable, and for $x \in [a, b]$, define
        \[
          G(x) = \int_{a}^{x} g.
        \]
        Then $G$ is continuous on $[a, b]$. If $g$ is continuous at some point $c \in [a, b]$,
        then $G$ is differentiable at $c$ and $G'(c) = g(c)$.
    \end{enumerate}
  \end{theorem}
\end{callout}

\begin{problem}
  Let $F(x) = \int_{0}^{x} t^{2} \dt$. Find $F'(x)$.
  \vspace{\baselineskip}

  $F'(x) = x^{2}$ by \thmref{thm:ftc}\ref{thm:ftc-ii}.
\end{problem}

\begin{problem}
  Let $G(x) = \int_{2}^{x} e^{t} \dt$. Compute $G'(x)$ and $G(2)$.
  \vspace{\baselineskip}

  $G'(x) = e^{x}, \quad G(2) = 0$ by \thmref{thm:ftc}\ref{thm:ftc-ii}.
\end{problem}

\begin{problem}
  Evaluate $\int_{1}^{4} 2x \dx$.
  \vspace{\baselineskip}

  $\int_{1}^{4} 2x \dx = \left. x^{2} \right\rvert_{1}^{4} = (4)^{2} - (1)^{2} = 15$ by \thmref{thm:ftc}\ref{thm:ftc-i}.
\end{problem}

\begin{problem}
  Find all continuous functions $f$ satisfying
  \[
    \int_{0}^{x} f(t) \dt = (f(x))^{2} + C
  \]
  for some constant $C$.
  \vspace{\baselineskip}

  Differentiating both sides using \thmref{thm:ftc}\ref{thm:ftc-ii} and
  \thmref{thm:chain-rule} gives
  \[
    f(x) = 2f(x)f'(x).
  \]
  Dividing by $f(x)$ (when nonzero) yields $f'(x) = \frac{1}{2}$,
  so $f(x) = \frac{x}{2} + A$. Evaluating at $x = 0$ gives $C = -A^{2}$.
  Therefore, $f(x) = \frac{x}{2} + A$ for any $A \in \R$, with $C = -A^{2}$.
\end{problem}

\begin{problem}
  \label{prob:leibniz-integral-rule}
  Let $F(x) = \int_{0}^{x} x \, f(t) \dt$. Find $F'(x)$. (The answer is not $x f(x)$;
  first rewrite the integral in an obvious way, then differentiate.)
  \vspace{\baselineskip}

  Since $x$ is constant with respect to $t$, we have
  $F(x) = x \int_{0}^{x} f(t) \dt$. By \thmref{thm:prod-rule-deriv}
  and \thmref{thm:ftc}\ref{thm:ftc-ii},
  \[
    F'(x) = \int_{0}^{x} f(t) \dt + x \, f(x).
  \]
  \begin{callout}[Leibniz's Integral Rule]
    \probref{prob:leibniz-integral-rule} is a special case of \textit{Leibniz's integral rule}
    (differentiation under the integral sign). In general, if
    $g(x, t)$ has a continuous partial derivative in $x$, then
    \[
      \frac{d}{dx} \int_{a(x)}^{b(x)} g(x, t) \dt
      = g(x, b(x)) \cdot b'(x) - g(x, a(x)) \cdot a'(x)
        + \int_{a(x)}^{b(x)} \frac{\partial g}{\partial x}(x, t) \dt.
    \]
    Here $g(x, t) = x f(t)$, $a(x) = 0$, and $b(x) = x$, so the
    boundary term gives $x f(x)$ and the partial derivative term gives
    $\int_{0}^{x} f(t) \dt$, recovering the same result.
  \end{callout}
\end{problem}

\begin{problem}
  Decide whether each statement is true or false, providing a short justification
  for each conclusion.
  \begin{enumerate}[label=(\alph*)]
    \item If $g = h'$ for some $h$ on $[a, b]$, then $g$ is continuous on $[a, b]$.

      False. If $h(x) = \abs{x}$ over $[-1, 1]$, then $g = h'$ is not continuous on 
      $[a,b]$.

    \item If $g$ is continuous on $[a, b]$, then $g = h'$ for some $h$ on $[a, b]$.

      True. Define $h(x) = \int_{a}^{x} g$. Since $g$ is continuous on $[a, b]$,
      \thmref{thm:ftc}\ref{thm:ftc-ii} implies $h'(x) = g(x)$ for all $x \in [a, b]$.

    \item If $H(x) = \int_{a}^{x} h$ is differentiable at $c \in [a, b]$, then $h$ is
      continuous at $c$.

      False. Let $h : [-1, 1] \to \R$ be defined by $h(x) = 1$ for $x \neq 0$ and
      $h(0) = 0$. Then $H(x) = \int_{-1}^{x} h = x + 1$, so $H'(0) = 1$ exists, but $h$
      is not continuous at $c = 0$.

  \end{enumerate}
\end{problem}

\begin{problem}
  The hypothesis in \thmref{thm:ftc}\ref{thm:ftc-i} that $F'(x) = f(x)$ for all
  $x \in [a, b]$ is slightly stronger than it needs to be. Carefully read the proof
  and state exactly what needs to be assumed with regard to the relationship between
  $f$ and $F$ for the proof to be valid.
  \vspace{\baselineskip}

  The proof applies \thmref{thm:mean-value-theorem} to $F$ on subintervals of a partition,
  which only requires $F$ to be continuous on $[a, b]$ and differentiable on $(a, b)$. Thus
  it suffices to assume $F'(x) = f(x)$ for all $x \in (a, b)$, with $F$ continuous on
  $[a, b]$.
\end{problem}

\begin{problem}
  Show that if $f : [a, b] \to \R$ is continuous and $\int_{a}^{x} f = 0$ for all
  $x \in [a, b]$, then $f(x) = 0$ everywhere on $[a, b]$. Provide an example to show
  that this conclusion does not follow if $f$ is not continuous.
  \vspace{\baselineskip}

  Define $F(x) = \int_{a}^{x} f = 0$ for all $x \in [a, b]$. Since $f$ is continuous,
  \thmref{thm:ftc}\ref{thm:ftc-ii} gives $F'(x) = f(x)$ for all $x \in [a, b]$. But
  $F \equiv 0$, so $f(x) = F'(x) = 0$ everywhere on $[a, b]$.

  For a counterexample without continuity, let $f(x) = 0$ for $x \neq c$ and $f(c) = 1$
  for some $c \in [a, b]$. Then $\int_{a}^{x} f = 0$ for all $x$, but $f(c) \neq 0$.
\end{problem}

\begin{problem}
  \label{prob:integration-by-parts}
  (Integration-by-parts).
  \begin{enumerate}[label=(\alph*)]
    \item Assume $h(x)$ and $k(x)$ have continuous derivatives on $[a, b]$ and derive
      the familiar integration-by-parts formula
      \[
        \int_{a}^{b} h(t)k'(t) \dt = h(b)k(b) - h(a)k(a) - \int_{a}^{b} h'(t)k(t) \dt.
      \]

      Define $F(x) = h(x)k(x)$. By \thmref{thm:prod-rule-deriv},
      $F'(x) = h'(x)k(x) + h(x)k'(x)$, which is continuous since $h', k, h, k'$ are all
      continuous. Applying \thmref{thm:ftc}\ref{thm:ftc-i},
      \[
        \int_{a}^{b} \bigl[h'(t)k(t) + h(t)k'(t)\bigr] \dt = F(b) - F(a) = h(b)k(b) - h(a)k(a).
      \]
      Rearranging gives the result.

    \item Explain how the result in Exercise 7.4.6 can be used to slightly weaken the
      hypothesis in part (a).

      Exercise 7.4.6 states that the product of integrable functions is integrable. Since
      $h$ and $k$ are differentiable and hence continuous, the products $h'k$ and $hk'$ are
      integrable whenever $h'$ and $k'$ are merely integrable. Thus $F' = h'k + hk'$ is
      integrable and \thmref{thm:ftc}\ref{thm:ftc-i} still applies, so we only need $h'$
      and $k'$ to be integrable rather than continuous.
  \end{enumerate}
\end{problem}

\begin{problem}
  Use part \ref{thm:ftc-ii} of \thmref{thm:ftc} to construct another proof of
  part \ref{thm:ftc-i} of \thmref{thm:ftc} under the stronger hypothesis that $f$
  is continuous. (To get started, set $G(x) = \int_{a}^{x} f$.)
  \begin{proof}
    Define $G(x) = \int_{a}^{x} f$. Since $f$ is continuous,
    \thmref{thm:ftc}\ref{thm:ftc-ii} gives $G'(x) = f(x)$ for all $x \in [a, b]$.
    But $F'(x) = f(x)$ as well, so $(F - G)'(x) = 0$ for all $x \in [a, b]$.

    Let $D = F - G$. For any $x \in [a, b]$, applying
    \thmref{thm:mean-value-theorem} to $D$ on $[a, x]$ yields some $c \in (a, x)$
    with
    \[
      D(x) - D(a) = D'(c)(x - a) = 0 \cdot (x - a) = 0,
    \]
    so $D(x) = D(a)$ for all $x$, i.e.\ $F(x) = G(x) + C$ for some constant $C$.
    Since $G(a) = 0$,
    \[
      F(b) - F(a) = (G(b) + C) - (G(a) + C) = G(b) = \int_{a}^{b} f.
    \]
  \end{proof}
\end{problem}

\begin{problem}
  \label{prob:natural-log-euler}
  (Natural Logarithm and Euler's Constant). Let
  \[
    L(x) = \int_{1}^{x} \frac{1}{t} \dt,
  \]
  where we consider only $x > 0$.
  \begin{enumerate}[label=(\alph*)]
    \item What is $L(1)$? Explain why $L$ is differentiable and find $L'(x)$.

      $L(1) = \int_{1}^{1} \frac{1}{t} \dt = 0$. Since $1/t$ is continuous for $x > 0$,
      \thmref{thm:ftc}\ref{thm:ftc-ii} implies $L$ is differentiable and $L'(x) = 1/x$.

    \item Show that $L(xy) = L(x) + L(y)$. (Think of $y$ as a constant and
      differentiate $g(x) = L(xy)$.) \label{prob:log-product-rule}

      Fix $y > 0$ and define $g(x) = L(xy)$. By \thmref{thm:ftc}\ref{thm:ftc-ii} and
      \thmref{thm:chain-rule}, $g'(x) = \frac{1}{xy} \cdot y = \frac{1}{x} = L'(x)$. So $(g - L)'(x) = 0$, and
      by \thmref{thm:mean-value-theorem}, $g(x) = L(x) + C$ for some constant $C$. Setting
      $x = 1$ gives $L(y) = L(1) + C = C$. Therefore $L(xy) = L(x) + L(y)$.

t a    \item Show $L(x/y) = L(x) - L(y)$.

      By \ref{prob:log-product-rule}, $L(x) = L\bigl(\frac{x}{y} \cdot y\bigr) = L(x/y) + L(y)$.
      Rearranging gives $L(x/y) = L(x) - L(y)$.

    \item Let
      \[
        \gamma_{n} = \left( 1 + \frac{1}{2} + \frac{1}{3} + \cdots + \frac{1}{n} \right) - L(n).
      \]
      Prove that $(\gamma_{n})$ converges. The constant $\gamma = \lim \gamma_{n}$ is called
      Euler's constant.

      \begin{proof}
        We show $(\gamma_{n})$ is decreasing and bounded below, then apply
        \thmref{thm:monotone-convergence}. First, observe
        \begin{align*}
          \gamma_{n} - \gamma_{n+1}
            &= -\frac{1}{n+1} + \left( \int_{1}^{n+1} \frac{1}{t} \dt - \int_{1}^{n} \frac{1}{t} \dt \right) \\
            &= \int_{n}^{n+1} \frac{1}{t} \dt - \frac{1}{n+1}.
        \end{align*}
        Since $1/t$ is strictly decreasing on $[n, n+1]$,
        \[
          \int_{n}^{n+1} \frac{1}{t} \dt > \int_{n}^{n+1} \frac{1}{n+1} \dt = \frac{1}{n+1},
        \]
        so $\gamma_{n} > \gamma_{n+1}$ for all $n$. Next, because $1/t$ is decreasing,
        each term $1/k$ satisfies
        \[
          \frac{1}{k} \geq \int_{k}^{k+1} \frac{1}{t} \dt
        \]
        for $k = 1, \ldots, n-1$. Summing gives
        $\sum_{k=1}^{n-1} \frac{1}{k} \geq \int_{1}^{n} \frac{1}{t} \dt$, so
        \[
          \gamma_{n} = \sum_{k=1}^{n} \frac{1}{k} - L(n) \geq \frac{1}{n} > 0.
        \]
        Therefore $(\gamma_{n})$ is decreasing and bounded below by $0$, and
        \thmref{thm:monotone-convergence} gives the result.
      \end{proof}

    \item Show how consideration of the sequence $\gamma_{2n} - \gamma_{n}$ leads to the
      interesting identity
      \[
        L(2) = 1 - \frac{1}{2} + \frac{1}{3} - \frac{1}{4} + \frac{1}{5} - \frac{1}{6} + \cdots.
      \]
  \end{enumerate}
\end{problem}

\begin{problem}
  \label{prob:total-variation}
  Given a function $f$ on $[a, b]$, define the \textit{total variation} of $f$ to be
  \[
    Vf = \sup \left\{ \sum_{k=1}^{n} \abs{f(x_{k}) - f(x_{k-1})} \right\},
  \]
  where the supremum is taken over all partitions $P$ of $[a, b]$.
  \begin{enumerate}[label=(\alph*)]
    \item If $f$ is continuously differentiable ($f'$ exists as a continuous function),
      use the Fundamental Theorem of Calculus to show $Vf \leq \int_{a}^{b} \abs{f'}$.

    \item Use the Mean Value Theorem to establish the reverse inequality and conclude
      that $Vf = \int_{a}^{b} \abs{f'}$.
  \end{enumerate}
\end{problem}
